\documentclass[10pt,a4paper]{amsart}
\usepackage{amsmath,amssymb,amsthm,graphicx}
% For commutative diagrams we use Xy-pic
\usepackage[all]{xy}

% We use 2cell for 2-commutative diagrams.
\xyoption{2cell}
% Theorem environments.
%
\theoremstyle{plain}
\newtheorem{theorem}[subsection]{Theorem}
\newtheorem{proposition}[subsection]{Proposition}
\newtheorem{lemma}[subsection]{Lemma}

\theoremstyle{definition}
\newtheorem{definition}[subsection]{Definition}
\newtheorem{example}[subsection]{Example}
\newtheorem{exercise}[subsection]{Exercise}
\newtheorem{situation}[subsection]{Situation}

\theoremstyle{remark}
\newtheorem{remark}[subsection]{Remark}
\newtheorem{remarks}[subsection]{Remarks}

\numberwithin{equation}{subsection}

% Macros
%
\def\lim{\mathop{\mathrm{lim}}\nolimits}
\def\colim{\mathop{\mathrm{colim}}\nolimits}
\def\Spec{\mathop{\mathrm{Spec}}}
\def\Hom{\mathop{\mathrm{Hom}}\nolimits}
\def\Ext{\mathop{\mathrm{Ext}}\nolimits}
\def\SheafHom{\mathop{\mathcal{H}\!\mathit{om}}\nolimits}
\def\SheafExt{\mathop{\mathcal{E}\!\mathit{xt}}\nolimits}
\def\Sch{\mathit{Sch}}
\def\Mor{\mathop{\mathrm{Mor}}\nolimits}
\def\Ob{\mathop{\mathrm{Ob}}\nolimits}
\def\Sh{\mathop{\mathit{Sh}}\nolimits}
\def\NL{\mathop{N\!L}\nolimits}
\def\CH{\mathop{\mathrm{CH}}\nolimits}
\def\proetale{{pro\text{-}\acute{e}tale}}
\def\etale{{\acute{e}tale}}
\def\QCoh{\mathit{QCoh}}
\def\Ker{\mathop{\mathrm{Ker}}}
\def\Im{\mathop{\mathrm{Im}}}
\def\Coker{\mathop{\mathrm{Coker}}}
\def\Coim{\mathop{\mathrm{Coim}}}

% Boxtimes
%
\DeclareMathSymbol{\boxtimes}{\mathbin}{AMSa}{"02}

%
% Macros for moduli stacks/spaces
%
\def\QCohstack{\mathcal{QC}\!\mathit{oh}}
\def\Cohstack{\mathcal{C}\!\mathit{oh}}
\def\Spacesstack{\mathcal{S}\!\mathit{paces}}
\def\Quotfunctor{\mathrm{Quot}}
\def\Hilbfunctor{\mathrm{Hilb}}
\def\Curvesstack{\mathcal{C}\!\mathit{urves}}
\def\Polarizedstack{\mathcal{P}\!\mathit{olarized}}
\def\Complexesstack{\mathcal{C}\!\mathit{omplexes}}
% \Pic is the operator that assigns to X its picard group, usage \Pic(X)
% \Picardstack_{X/B} denotes the Picard stack of X over B
% \Picardfunctor_{X/B} denotes the Picard functor of X over B
\def\Pic{\mathop{\mathrm{Pic}}\nolimits}
\def\Picardstack{\mathcal{P}\!\mathit{ic}}
\def\Picardfunctor{\mathrm{Pic}}
\def\Deformationcategory{\mathcal{D}\!\mathit{ef}}
\newenvironment{reference}{}{}
\newenvironment{slogan}{}{}
\newenvironment{history}{}{}
\usepackage{cleveref}
\crefname{originallemma}{original lemma}{original lemmas}
\Crefname{originallemma}{Original lemma}{Original lemmas}
\makeatletter
\newlabel{lemma-divisor-special-fiber}{{0C5Z}{source}{}{sourcepointer-1}{}}
\newlabel{lemma-divisor-special-fiber@cref}{{[originallemma][0][]0C5Z}{[1][0][]source}{}{}{}}
\makeatother
\crefname{originallemma}{original lemma}{original lemmas}
\Crefname{originallemma}{Original lemma}{Original lemmas}
\makeatletter
\newlabel{divisors-lemma-extend-invertible-module}{{0BD9}{source}{}{sourcepointer-2}{}}
\newlabel{divisors-lemma-extend-invertible-module@cref}{{[originallemma][0][]0BD9}{[1][0][]source}{}{}{}}
\makeatother
\crefname{originaldefinition}{original definition}{original definitions}
\Crefname{originaldefinition}{Original definition}{Original definitions}
\makeatletter
\newlabel{divisors-definition-divisor-invertible-sheaf}{{0BE6}{source}{}{sourcepointer-3}{}}
\newlabel{divisors-definition-divisor-invertible-sheaf@cref}{{[originaldefinition][0][]0BE6}{[1][0][]source}{}{}{}}
\makeatother
\crefname{originallemma}{original lemma}{original lemmas}
\Crefname{originallemma}{Original lemma}{Original lemmas}
\makeatletter
\newlabel{divisors-lemma-normal-c1-injective}{{0BE8}{source}{}{sourcepointer-4}{}}
\newlabel{divisors-lemma-normal-c1-injective@cref}{{[originallemma][0][]0BE8}{[1][0][]source}{}{}{}}
\makeatother
\title{The Picard group of a regular model sits in an exact sequence determined by the multiplicities of the special fibre}
\author{The Stacks Project Authors}
\date{stacks-project:models.tex:local-snapshot-20261012; GNU FDL 1.2 or later}
\begin{document}
\maketitle
\noindent\textit{Editorial note.} The counted unit is the Stacks Project Lemma lemma-regular-model-pic at source lines 4649-4658 together with its proof at source lines 4660-4681, a straight adjacent pairing in the chapter file. The chapter is the upstream file models.tex, which is the file name recorded in the source record; the pool directory carries the longer name stacks-models and is not used as the chapter name anywhere in this card.\par\medskip
\noindent\textit{Editorial note.} The retained statement opens with the words In Situation followed by a reference to situation-regular-model, so the situation environment that defines the setting is retained as a context span at source lines 4628-4637. It is context, not a counted unit, and it does not overlap the counted interval.\par\medskip
\noindent\textit{Editorial note.} This chapter's candidate index records the statement body and the proof body without their closing markers, so the retained spans used here are the full environments, exactly as in the earlier cards of this chapter.\par\medskip
\noindent\textit{Editorial note.} The retained proof quotes three results of the Divisors chapter, for the extension of invertible sheaves, for the divisor attached to a meromorphic section and for the injectivity of the divisor class map on a normal scheme; all three are declared as non-load-bearing navigation pointers with their stable Stacks tags looked up in the shared tag index. The retained situation additionally names one lemma of this chapter, also declared as a pointer. No citation occurs in the retained spans, so no bibliography entry is supplied.\par\medskip
\noindent\textit{Editorial note.} The retained preamble is the chapter preamble's Xy-pic setup at lines 14-18, the theorem declarations and mathematical macros at lines 154-217, and the three no-op wrapper environments.\par\medskip
\noindent\textit{Allowed standard prerequisites (editorial).}\begin{itemize}
\item Divisors, invertible sheaves and the restriction map on the Picard group. For an integral normal scheme the divisors of meromorphic sections of invertible sheaves give the usual presentation of the Picard group, restriction of an invertible sheaf to an open subscheme is surjective for the schemes at hand here, and the divisor of a rational function determines that function up to a unit. The author's own definitions and the three lemmas of the Divisors chapter quoted inside the proof are the allowed inputs.. The counted exact sequence is the tool the whole section uses to move between divisors on the model, divisors on the curve and the components of the special fibre.
\end{itemize}\medskip
\ifdefined\hypertarget\hypertarget{sourcepointer-1}{}\fi
\noindent\textit{Source reference.} Original lemma 0C5Z: focused-seeds/AGP/sources/stacks-snapshot-20261011/tags.txt, source lines 15279--15279. Navigation pointer inside the retained situation, which names it when it writes the special fibre as a sum of its components with multiplicities; the counted exact sequence uses those multiplicities as data. The number is the stable Stacks tag of the target, looked up in the shared tag index.\par\medskip
\ifdefined\hypertarget\hypertarget{sourcepointer-2}{}\fi
\noindent\textit{Source reference.} Original lemma 0BD9: focused-seeds/AGP/sources/stacks-snapshot-20261011/tags.txt, source lines 14307--14307. Allowed input of the Divisors chapter quoted by the retained proof for the surjectivity of restriction of invertible sheaves; the pointer is registered so the reference resolves and the cited statement is not reproduced inside the span. The number is the stable Stacks tag of the target, looked up in the shared tag index.\par\medskip
\ifdefined\hypertarget\hypertarget{sourcepointer-3}{}\fi
\noindent\textit{Source reference.} Original definition 0BE6: focused-seeds/AGP/sources/stacks-snapshot-20261011/tags.txt, source lines 14339--14339. Allowed input of the Divisors chapter quoted by the retained proof for the divisor of a meromorphic section of an invertible sheaf; the pointer is registered so the reference resolves and the cited definition is not reproduced inside the span. The number is the stable Stacks tag of the target, looked up in the shared tag index.\par\medskip
\ifdefined\hypertarget\hypertarget{sourcepointer-4}{}\fi
\noindent\textit{Source reference.} Original lemma 0BE8: focused-seeds/AGP/sources/stacks-snapshot-20261011/tags.txt, source lines 14341--14341. Allowed input of the Divisors chapter quoted by the retained proof for the injectivity of the divisor class map on a normal scheme; the pointer is registered so the reference resolves and the cited statement is not reproduced inside the span. The number is the stable Stacks tag of the target, looked up in the shared tag index.\par\medskip
\begin{situation}
\label{situation-regular-model}
Let $R$ be a discrete valuation ring with fraction field $K$,
residue field $k$, and uniformizer $\pi$.
Let $C$ be a smooth projective curve over $K$ with $H^0(C, \mathcal{O}_C) = K$.
Let $X$ be a regular proper model of $C$.
Let $C_1, \ldots, C_n$ be the irreducible components of the special
fibre $X_k$. Write $X_k = \sum m_i C_i$ as in
Lemma \ref{lemma-divisor-special-fiber}.
\end{situation}

\begin{lemma}
\label{lemma-regular-model-pic}
In Situation \ref{situation-regular-model} there is an exact sequence
$$
0 \to \mathbf{Z} \to \mathbf{Z}^{\oplus n} \to
\Pic(X) \to \Pic(C) \to 0
$$
where the first map sends $1$ to $(m_1, \ldots, m_n)$ and the second
maps sends the $i$th basis vector to $\mathcal{O}_X(C_i)$.
\end{lemma}

\begin{proof}
Observe that $C \subset X$ is an open subscheme. The restriction
map $\Pic(X) \to \Pic(C)$ is surjective by
Divisors, Lemma \ref{divisors-lemma-extend-invertible-module}.
Let $\mathcal{L}$ be an invertible $\mathcal{O}_X$-module
such that there is an isomorphism $s : \mathcal{O}_C \to \mathcal{L}|_C$.
Then $s$ is a regular meromorphic section of $\mathcal{L}$
and we see that $\text{div}_\mathcal{L}(s) = \sum a_i C_i$
for some $a_i \in \mathbf{Z}$
(Divisors, Definition \ref{divisors-definition-divisor-invertible-sheaf}).
By Divisors, Lemma \ref{divisors-lemma-normal-c1-injective}
(and the fact that $X$ is normal)
we conclude that $\mathcal{L} = \mathcal{O}_X(\sum a_iC_i)$.
Finally, suppose that $\mathcal{O}_X(\sum a_i C_i) \cong \mathcal{O}_X$.
Then there exists an element $g$ of the function field of $X$
with $\text{div}_X(g) = \sum a_i C_i$. In particular the rational
function $g$ has no zeros or poles on the generic fibre $C$ of $X$.
Since $C$ is a normal scheme this implies $g \in H^0(C, \mathcal{O}_C) = K$.
Thus $g = \pi^a u$ for some $a \in \mathbf{Z}$ and $u \in R^*$.
We conclude that $\text{div}_X(g) = a \sum m_i C_i$ and the proof
is complete.
\end{proof}

\end{document}
